\section*{MEASURES ON SEQUENCE SPACES.}
\phantomsection
\addcontentsline{toc}{section}{MEASURES ON SEQUENCE SPACES.}

One of the most developed branches of classical Probability Theory is that of limit theorems (the law of large numbers, tendency towards the Gauss-Laplace law, \ldots); it is a case of a deepening of the notion of statistical regularity manifested by phenomena bringing into play very large populations. The correct mathematical formulation of these problems requires the introduction of measures on the space of sequences; these spaces, which constitute the most obvious generalisation of spaces of finite dimension, are the favourite subject of the researches on ``General Analysis'' undertaken around 1920 by Fréchet, Lévy, Lusin, \ldots{} It is in any case not fortuitous that Khintchine and Kolmogoroff, the creators of the new methods in Probability Theory, are both disciples of Lusin, and that Lévy very soon turned to probabilistic problems: these constituted the touch stone of the new methods.

The first implicit occurrence of a measure on a sequence space appears in the work devoted by E. Borel in 1909 to countable probabilities {[}32 b{]}. A very original idea of Borel consisted of the application of probabilistic results that he had just obtained to the proof of properties possessed by the decimal expansion of nearly every real number contained between 0 and 1. This application rests on the following fundamental remark: let us define every real number contained between 0 and 1 by the sequence of digits in its expansion in a given base \(q\) ( \(q \geq 2\) ); if one takes at random successively the different digits from a number \(x\) , independently the ones from the others and with an equal probability \(1/q\) for \(0,1,\ldots,q-1\) , the probability that \(x\) is to be found in an interval of [0,1] is equal to the length of this interval.

In 1923, Steinhaus {[}293{]} establishes these results rigorously and describes the precise mathematical model of the unbounded sequence of selections of the kind considered by Borel: let us take \(q = 2\) to simplify matters and denote by \(I\) the set with two elements \(\{0,1\}\) ; we furnish \(I\) with the measure \(\mu\) defined by \(\mu(0) = \mu(1) = \frac{1}{2}\) ; the elements of the product space \(I^{\mathbf{N}}\) are the sequences \(\varepsilon = (\varepsilon(n))_{n \in \mathbf{N}}\) of numbers equal to 0 or 1 and the map \(\phi : \varepsilon \mapsto \sum_{n \geq 0} \varepsilon(n). 2^{-n-1}\) is, up to a countable set, a bijection of \(I^{\mathbf{N}}\) on the interval \([0,1]\) ; what is more, \(\phi^{-1}\) transforms the Lebesgue measure on \([0,1]\) into the measure \(P\) on \(I^{\mathbf{N}}\) the product of the measures \(\mu\) on each of the factors. In fact Steinhaus does not have at his disposal a construction for product measures;

he uses the existence of the quasi-bijection \(\phi\) to construct the measure P on \(I^{N}\) starting from the Lebesgue measure on [0,1], then he gives an axiomatic characterisation of P. The isomorphism thus obtained allows the translation of the language of probability into that of measure and to apply the known theorems about the Lebesgue integral.

In the same work, Steinhaus considers the random series \(\sum_{n\geq 0}\sigma_n.a_n\) , where the signs \(\sigma_{n} = \pm 1\) are chosen at random independently the ones from the others and with the same probability \(\frac{1}{2}\) ; between 1928 and 1935, he studies numerous other random series. On their side, Paley, Wiener and Zygmund consider random Fourier series \(^{1}\) of the form \(\sum_{-\infty}^{\infty}a_n\exp (2\pi i(nt + \varPhi_n))\) ; the ``amplitudes'' \(a_{n}\) are fixed, and the ``phases'' \(\varPhi_{n}\) are independent random variables uniformly spread over [0, 1]. If the analytic difficulties vary enormously from problem to problem, the translation in terms of measure theory is the same in all cases and represents an extension of the case treated by Borel and Steinhaus; it is a matter of constructing a measure on \(\mathbf{R}^{\mathbf{N}}\) , the product of a family of measures all identical to one particular positive measure \(\mu\) of mass 1 on \(\mathbf{R}\) ; for example, the preceding random Fourier series correspond to the case where \(\mu\) is the Lebesgue measure on [0, 1].

To construct such product measures, two methods can be used. The first is a direct method, set up for the first time by Daniell {[}76 b{]} in 1918; it is found again in 1934 by Jessen {[}173{]} who will make a detailed study of the case where \(\mu\) is the Lebesgue measure on \([0,1]\) . The second method is the search for tricks analogous to those of Steinhaus in order to reduce to the case of the Lebesgue measure on \([0,1]\) ; this way of proceeding had the advantage of convenience as long as there is not at our disposal a complete account of the general theory of measure, for it allows the use without new proofs of the Lebesgue theorems. \(^{2}\)

